\documentclass[10pt,a4paper]{article}

\usepackage[utf8]{inputenc}
\usepackage[T1]{fontenc}
\usepackage[ngerman]{babel}
\usepackage[margin=2cm]{geometry}
\usepackage{graphicx}
\usepackage{listings}
\usepackage{hyperref}

% ---------------------------------------------------------------------------
% Bilder liegen in img/. Damit das Dokument aus jedem Arbeitsverzeichnis
% uebersetzt (aus latex/ selbst, aus dem Repo-Wurzelverzeichnis oder aus dem
% Paketverzeichnis), werden hier mehrere Suchpfade angegeben.
% Hinweis: Unterstriche in Dateinamen vermeiden, LaTeX stolpert darueber.
% ---------------------------------------------------------------------------
\graphicspath{%
  {img/}%
  {latex/img/}%
  {docs/latex/img/}%
  {src/packages/obstacle_detection/docs/latex/img/}%
  {../img/}%
  {../latex/img/}%
}

% Abbildung mit vorhandener Datei.
% #1 = Dateiname (ohne Verzeichnis), #2 = Breite, #3 = Bildunterschrift
\newcommand{\bild}[3]{%
  \par\medskip\noindent\begin{center}
  \includegraphics[width=#2\linewidth,height=0.8\textheight,keepaspectratio]{#1}\\[0.5em]
  \begin{minipage}{0.9\linewidth}\footnotesize #3\end{minipage}
  \end{center}\medskip}

\lstset{basicstyle=\ttfamily\footnotesize, frame=single, framesep=4pt,
        columns=fullflexible, keepspaces=true, xleftmargin=1em}

\title{\textbf{obstacle\_detection}\\[0.3em]
       \large Enten erkennen und umfahren für die Duckiebot Challenge}
\author{Nico Sander}
\date{\today}
\usepackage{titlesec}
\titlespacing*{\section}{0pt}{1.2ex plus .5ex}{0.8ex}
\titlespacing*{\subsection}{0pt}{1ex plus .3ex}{0.5ex}
\setlength{\parskip}{0.4ex}

\begin{document}
\maketitle

\section{Ausgangslage und Aufgabenstellung}

\subsection{Was gegeben ist}

Gegeben ist ein Kurs aus Klebeband auf schwarzen Matten. Eine zweispurige
Fahrbahn, getrennt durch eine gestrichelte gelbe Mittellinie und außen begrenzt
durch durchgezogene weiße Linien, mündet in einen Wendeplatz. Auf dem
Wendeplatz stehen Gummienten.

\bild{kurs.jpg}{0.34}{Der Kurs. Unten die Startposition auf der rechten Spur,
oben der Wendeplatz mit den Enten. Innerhalb des Wendeplatzes laufen die
Fahrbahnmarkierungen nicht weiter.}

Der Duckiebot bringt für die Aufgabe eine nach vorn gerichtete Kamera,
Radencoder und einen Abstandssensor mit.

\subsection{Zu lösendes Problem}

Die Aufgabe lautet: eine Runde durch den Wendeplatz fahren, ohne eine Ente
umzufahren.

Gewertet wird nicht die Zeit. Die Geschwindigkeit spielt keine Rolle, und es
gibt keinen Vorteil darin, schnell zu sein. Bewertet wird stattdessen, mit wie
vielen Enten auf dem Wendeplatz die Umsetzung noch zurechtkommt. Wie die Enten
stehen, ist vorher nicht bekannt.

Aus dem letzten Punkt folgt die eigentliche Schwierigkeit: eine feste, vorab
geplante Route gibt es nicht, denn sie müsste für eine Anordnung ausgelegt
werden, die zur Laufzeit eine andere sein wird. Der Bot muss also während der
Fahrt entscheiden, wo er durchpasst.

Daraus ergeben sich vier Teilprobleme:

\begin{itemize}
  \itemsep0.2ex
  \item der Fahrbahn folgen und erkennen, wo sie verläuft,
  \item die Enten finden und abschätzen, wie weit sie entfernt sind,
  \item entscheiden, ob ausgewichen werden muss und wohin,
  \item dabei nicht dauerhaft stehen bleiben, denn eine nicht beendete Runde
        zählt nicht.
\end{itemize}

\section{Systemüberblick und Kernkomponenten}

Das Paket besteht aus vier Nodes. Drei davon bilden zwei Schichten: die
\emph{Wahrnehmung} wertet das Kamerabild aus und liefert daraus die Lage der
Fahrbahn und die Positionen der Enten, die \emph{Fahrentscheidung} verrechnet
beides zu einem Fahrbefehl. Der vierte Node ist ein Dashboard, das rein der
Fehlersuche dient und am Fahrverhalten nichts ändert.

Von den vorhandenen Sensoren wird nur die Kamera genutzt; Encoder und
Abstandssensor bleiben ungenutzt. Der Regler entscheidet jeweils aus dem
aktuellen Kamerabild neu und führt keine Karte der Umgebung mit.

\bild{nodegraph.png}{0.42}{Der ROS Node Graph, auf den Hauptdatenfluss
reduziert. Kästen sind Nodes, Pfeile sind Topics, gestrichelt sind Kamera und
Radansteuerung als externe Beteiligte. Blau ist die Wahrnehmung, rot die
Fahrentscheidung, violett die Anzeige. Die Debugbilder ans Dashboard sind der
Übersicht halber weggelassen.}

\subsection{Wahrnehmung}\label{sec:wahrnehmung}

\paragraph{\texttt{detect\_lane\_node.py}} Segmentiert das Kamerabild in weiße
und gelbe Fahrbahnmarkierungen und leitet daraus den Querabstandsfehler zur
Spurmitte und die Spaltenpositionen der beiden Linien ab. Das U-Net dahinter ist
dasselbe wie in den Challenges 1 und 2 und dort beschrieben; neu sind die
folgenden drei Punkte.

Für diese Challenge kam eine Schutzmaßnahme dazu: bevor die Linienpositionen
berechnet werden, werden die aktuellen Entenboxen aus der Segmentierungsmaske
entfernt. Eine gelbe Ente auf der Fahrbahn sieht dem gelben Mittelstreifen
sonst ähnlich genug, um die Spurlage zu verfälschen, und zwar genau in dem
Moment, in dem der Bot auf sie zufährt.

Zwei weitere Signale entstehen hier und werden zusammen mit den
Linienpositionen verschickt. Verteilen sich die Pixel einer Linie über einen
großen Teil der Bildbreite, statt einer Spalte zu folgen, so liegt sie quer im
Bild und wird als frontal angefahren gemeldet: sie ist dann keine
Spurbegrenzung, der man folgen kann, sondern eine Wand. Und erscheinen die
beiden Linienfarben in vertauschter Reihenfolge, ist der Bot verkehrt herum in
der Spur unterwegs.

\paragraph{\texttt{detect\_obstacle\_node.py}} Erkennt Enten mit einem selbst
trainierten YOLOv11 Modell und veröffentlicht alle Bounding Boxen als JSON. Der
Node trifft keine Fahrentscheidung, er liefert nur strukturierte
Hindernisinformation.

Wie \texttt{detect\_lane\_node} schützt er sich mit einer
Wiedereintrittssperre: liegt die Inferenz hinter dem Kamerabildstrom zurück,
wird das neue Bild verworfen statt eingereiht. Ein verworfenes Bild ist besser
als ein veraltetes, denn gehandelt wird auf der aktuellen Pose.

\subsection{Fahrentscheidung}

\paragraph{\texttt{control\_lane\_node.py}} Der zentrale Node und das Herz des
Pakets. Er ist als Regler nach dem Prinzip \glqq Folge der Lücke\grqq{} mit
einem kleinen Zustandsautomaten aufgebaut. Das Spurfolgen liefert dabei die
Route einschließlich der Wende; die Ausweichschicht verändert nur die Lenkung
und klemmt sie hart in den Korridor zwischen den sichtbaren Linien.

Aufgebaut ist er als Zustandsautomat nach dem EVA Prinzip: Sensordaten hinein,
Zustand bestimmen, genau ein Fahrbefehl hinaus. Weil der Fahrbefehl an einer
einzigen Stelle entsteht, lässt sich dort garantieren, dass er immer aktiv ist.
Gibt es keinen Weg nach vorn, dreht der Bot sich also auf der Stelle und sucht
eine Lücke, statt anzuhalten. Findet er längere Zeit keine, lockert er
schrittweise seine Anforderungen, bis eine Lücke als befahrbar gilt.

Die eigentliche Fahrlogik steckt in einer eigenen Klasse. Der Node sammelt die
eingehenden Nachrichten ein und ruft sie zehnmal pro Sekunde auf; sie liefert
daraus Geschwindigkeit und Drehrate.

\paragraph{Startsperre} Der Regler fährt erst los, wenn alle drei Eingangsdaten
mindestens einmal eingetroffen sind, denn ein noch stummes Topic ist von
\glqq keine Ente voraus\grqq{} nicht zu unterscheiden.

\subsection{Bedienung und Visualisierung}

\paragraph{\texttt{dashboard\_node.py}} Fasst die Debugbilder in einer Ansicht
zusammen: Kamerabild mit erkannten Linien, die weiße und die gelbe Maske sowie
ein Overlay der Planung. Im Overlay sind die durch Enten gesperrten Bereiche
rot, die freien grün, die gewählte Lücke cyan umrandet und die aktuelle
Zielspalte blau. Damit ist auf einen Blick zu sehen, warum der Regler gerade so
entscheidet, wie er entscheidet.

\paragraph{Trockenlauf} Die Startdatei kennt ein Argument, das den gesamten
Stapel startet und planen lässt, aber den Fahrbefehl nicht veröffentlicht. Der
Bot steht dabei still, während das Dashboard vollständig gefüllt bleibt. Das
ist die Betriebsart, in der die Parameter abgestimmt werden.

\section{Zustands- und Ablaufdiagramme}

\subsection{Fahrzustände}

Der Zustandsautomat ist bewusst klein gehalten. Alle Übergänge liegen in einer
einzigen Funktion, getrieben von skalaren Signalen und Mindestverweilzeiten,
sodass es keine zweite Stelle gibt, die dieselbe Entscheidung anders treffen
könnte.

\bild{fahrzustaende.png}{0.38}{Die Fahrzustände des
\texttt{control\_lane\_node}. \texttt{WAITING} ist die Startsperre und der
einzige Zustand, in dem der Bot steht. \texttt{ESCAPE\_ROTATE} fährt nicht,
dreht aber immer, womit die Invariante gewahrt bleibt. In diesen Zustand führen
außer einer versperrten Fahrbahn auch die beiden Sonderfälle aus
Abschnitt~\ref{sec:wahrnehmung}: eine quer angefahrene Markierung und eine
verkehrt herum befahrene Spur. Beide werden wie ein Hindernis behandelt, weil
Drehen in beiden Fällen die richtige Antwort ist und die Lage zugleich auflöst.}

\subsection{Lückenplanung pro Takt}

In jedem Takt wird aus den beiden Sensorbildern eine einzige Zielspalte
berechnet. Die sichtbaren Linien sind dabei harte Grenzen: der Zielpunkt wird
immer in den Korridor geklemmt, sodass keine Markierung überfahren wird.

Hat der Bot eine Lücke gewählt, bleibt er für eine Mindestzeit dabei, auch wenn
sie zwischenzeitlich zu schmal misst: die Boxen wachsen beim Heranfahren, sodass
jede Lücke zwangsläufig enger wird und das Manöver sonst auf halbem Weg
abbrechen würde. Nur eine wirklich nahe Ente direkt voraus hebt diese Bindung
sofort auf.

Bevor die freien Intervalle bestimmt werden, wird jede Entenbox verbreitert,
und zwar umso stärker, je näher die Ente steht. Eine feste Verbreiterung im
Bild wäre nur in genau einer Entfernung richtig:
die scheinbare Breite eines Objekts fällt mit dem Kehrwert der Entfernung, und
der Bot selbst erscheint aus wenigen Zentimetern um ein Vielfaches breiter im
Bild als aus einem halben Meter. Sicherheitsabstände sind deshalb in
Zentimetern angegeben; eine auf dem Kurs gemessene Kennlinie übersetzt die
Unterkante einer Box in eine Entfernung und diese erst im Takt in Bildbreite.

\bild{luecken.png}{1.0}{Die Lückenplanung. Aus den freien Intervallen wird das
breiteste gewählt und das Ziel hineingeklemmt. Eine Hysterese dämpft das
Umschalten zwischen zwei ähnlich breiten Lücken, damit der Bot nicht zwischen
links und rechts pendelt.}

\section{Konfiguration}

Die abgestimmten Größen liegen in \texttt{config/} als JSON, getrennt nach
Regler und Detektor. Der Regler kann seine Parameter zusätzlich zur Laufzeit
entgegennehmen, was das Abstimmen auf dem Bot deutlich verkürzt.

Zum Abstimmen gehört außerdem ein Werkzeug, das eine aufgezeichnete Fahrt mit
geänderten Parametern erneut durchrechnet. Auf dem Roboter liefert jede Sitzung
nur eine verrauschte Stichprobe; dieselbe Aufzeichnung erneut zu rechnen, macht
daraus einen kontrollierten Vergleich.

Am stärksten auf das Verhalten wirken diese fünf Größen:

\begin{itemize}
  \itemsep0.2ex
  \item \texttt{safety\_margin\_cm} (2\,cm): Abstand, den der Bot zusätzlich zu
        seiner halben Breite um jede Ente herum freihält. Der größte Hebel
        dafür, wie eng er an einer Ente vorbeifährt.
  \item \texttt{gap\_min\_cm} (2,4\,cm): wie viel breiter als dieser
        Mindestabstand eine Lücke sein muss, damit sie als befahrbar zählt.
  \item \texttt{front\_block\_ymax} (0,8): ab welcher Nähe eine Ente direkt
        voraus als Sperre gilt und der Bot dreht, statt weiterzufahren.
  \item \texttt{avoid\_min\_dwell} (1,6\,s): wie lange der Bot an einer einmal
        gewählten Lücke festhält.
  \item \texttt{escape\_relax\_after} (1,5\,s): nach welcher erfolglosen Drehzeit
        er beginnt, Mindestbreite und Sicherheitsabstand zu lockern.
\end{itemize}

\end{document}
